\section{Related Work}

\paragraph{Text Embedding Models.}
Text embedding models encode variable-length text into fixed-dimensional vectors for retrieval, clustering, and semantic textual similarity. Sentence-BERT made transformer encoders practical for sentence-level similarity search~\citep{reimers2019sentencebert}, SimCSE showed that contrastive learning with dropout noise or entailment supervision produces strong sentence embeddings~\citep{gao2021simcse}, and evaluation suites such as BEIR and MTEB shifted the field toward broad multi-task assessment~\citep{thakur2021beir,muennighoff2023mteb}. Recent systems scale this recipe with larger foundation models and specialized retrieval objectives: BGE-M3 combines multilingual, multi-functionality, and multi-granularity training with self-distillation across retrieval modes~\citep{chen2024bgem3}, and Qwen3 Embedding scales embedding and reranking models from Qwen3~\citep{zhang2025qwen3embedding}. This line supplies stronger teachers but does not specify how a compact student should inherit their geometry. RIPPLE treats such models as black-box teachers, requiring only their final embeddings.

\paragraph{Distillation for Language and Embedding Models.}
Knowledge distillation transfers behavior from a large teacher to a smaller student~\citep{hinton2015distilling}. In pretrained language models, DistilBERT, TinyBERT, and MiniLM compress transformer encoders through combinations of logits, hidden states, attention distributions, and self-attention relations~\citep{sanh2019distilbert,jiao2019tinybert,wang2020minilm}. Recent LLM distillation extends this toolkit with on-policy training, contrastive objectives, sparse-logit acceleration, and score-matching formulations~\citep{agarwal2024onpolicy,ko2025distillm2,anshumann2025sparse,kim2025csd,xu2024survey}, and cross-tokenizer methods such as DSKD and CDM address vocabulary mismatch~\citep{zhang2024dskd,chen2025cdm}.

Embedding-specific distillation adapts these ideas to retrieval-oriented representations. EmbedDistill aligns student and teacher retrieval geometry, including query--document embedding structure, for dense retrievers~\citep{kim2023embeddistill}. Gecko distills LLM-generated and LLM-relabelled retrieval data into compact embedding models~\citep{lee2024gecko}; Jasper and Stella combine multi-stage embedding distillation with Matryoshka Representation Learning~\citep{zhang2024jasperstella}; LEAF aligns compact students to teacher representations and supports asymmetric retrieval~\citep{vujanic2025leaf}; and Jina Embeddings v5 pairs task-targeted distillation with contrastive training~\citep{akram2026jinaembeddingsv5}. These methods establish distillation as a central recipe for embedding development. RIPPLE is complementary: rather than introducing teacher internals or task labels, it converts cached final embeddings into manifold diffusion targets.

\paragraph{Relational and Local Supervision.}
This is the family RIPPLE is positioned against. Pointwise embedding distillation matches a projected student embedding to the corresponding teacher embedding; it is architecture-agnostic but asks a compact student to reproduce absolute coordinates even when its capacity and dimension differ. Contrastive sentence-embedding distillation improves on this by transferring contrastive behavior and pairwise preferences~\citep{gao2021distilcse,xu2023distillcse}, and relational knowledge distillation argues that students should preserve relations among examples rather than individual outputs~\citep{park2019rkd}. In embedding training this becomes mini-batch similarity matching: teacher and student pairwise similarity matrices are compared over the examples that happen to co-occur in a stochastic batch. Such objectives reflect retrieval better than coordinate regression, but their supervision is bounded by the sampled batch, and Section~\ref{sec:theory} shows the bound is sharper than it appears---neighbor coverage accrues at rate $O(B/N)$ per epoch, and the normalized in-batch target is biased, not merely noisy.

A second line supplies intermediate supervision when teacher internals or student layers are available, including CKA-style representation matching and optimal-transport alignment~\citep{kornblith2019similarity,alqahtani2021ot}. Closest among these is teacher-anchored layer alignment~\citep{dao2026talas}, our strongest baseline, which targets internal transfer under capacity mismatch. RIPPLE requires no such access: it uses only the neighborhood geometry induced by final teacher embeddings.

\paragraph{Graph, Manifold, and Diffusion Geometry.}
Graph-based manifold learning studies how neighborhood graphs reveal structure in high-dimensional data. Diffusion Maps use Markov diffusion on a graph to expose multi-scale connectivity~\citep{coifman2006diffusionmaps}, and more recently HeatGeo relates heat diffusion to geodesic distance through Varadhan's formula, fitting an embedding that preserves a heat-geodesic dissimilarity~\citep{huguet2023heatgeo}. Knowledge distillation on graphs studies how structural information transfers between graph models or from graph teachers to simpler students~\citep{yang2023graphkd_survey,zhang2022glnn}.

RIPPLE draws on this view but inverts what diffusion is used for. Manifold-learning methods fit \emph{coordinates} for a fixed point set so that a diffusion-derived distance is preserved; RIPPLE uses diffusion rows as \emph{targets} for a parametric encoder, so what transfers is a conditional distribution over neighbors rather than a distance to be embedded, and the result applies to text the graph never contained. The graph is also not an observed input graph: it is induced by a strong text embedding teacher, which is what makes its diffusion distributions a supervision signal rather than an output.
